\chapter{Literature Review}
\label{ch:literature}

This chapter surveys the key intellectual traditions this thesis draws upon:
optimal growth theory under uncertainty, Bayesian inference in sequential
settings, multi-armed bandit algorithms, Hidden Markov Models for financial
regime detection, and portfolio insurance mechanisms. Where relevant, we
identify the specific limitations of each body of work that motivate our
proposed extensions.

%---------------------------------------------------------------
\section{The Kelly Criterion and Optimal Growth Theory}
%---------------------------------------------------------------

The foundational result in optimal betting theory is due to
\textcite{kelly1956new}, who considered a gambler receiving a noisy binary
signal and asked: what fraction of wealth should be wagered each round to
maximize long-run wealth? Kelly showed that the answer is to maximize the
expected logarithm of wealth, yielding the celebrated formula
\[
    f^* = \frac{bp - q}{b},
\]
where $p$ is the probability of winning, $q = 1-p$, and $b$ is the net odds.
The key insight is that log-utility is the unique utility function consistent
with asymptotically maximal wealth accumulation---any other utility leads to
eventual underperformance relative to Kelly betting in the long run.

The theoretical optimality of Kelly was made precise by \textcite{breiman1961},
who proved that the Kelly bettor almost surely accumulates more wealth than any
other essentially different strategy as the horizon $T \to \infty$. This
result, a direct consequence of the Strong Law of Large Numbers (SLLN) for
log-returns, forms the theoretical backbone of Chapter~\ref{ch:foundations}.

A major intellectual challenge to Kelly came from \textcite{samuelson1979},
who argued forcefully that log-utility is not intrinsically special and that
risk-averse investors with different utility functions should not follow Kelly.
The ``Samuelson--Kelly debate'' remains live in the portfolio management
literature. This thesis sides with the empirical tradition: we use Kelly as a
theoretical benchmark whose optimality holds under its stated assumptions, and
we study what happens when those assumptions---stationarity, correct probability
estimates---are violated.

For continuous return distributions, the Kelly fraction generalizes to
$f^* = \mu / \sigma^2$ when returns are drawn from $\Normal(\mu, \sigma^2)$,
a result obtained by maximizing $\E[\log(1 + fR)]$ for $R \sim \Normal(\mu,
\sigma^2)$ with respect to $f$ (see, e.g., \textcite{thorp2008}). This
continuous-distribution form is used throughout our simulation framework.

%---------------------------------------------------------------
\section{Bayesian Updating and the Sticky Prior Problem}
%---------------------------------------------------------------

Bayesian sequential decision-making treats parameter uncertainty as a
probability distribution over hypotheses that is updated as data arrives.
The canonical treatment for binary outcomes uses a Beta-Binomial conjugate
model: a prior $\text{Beta}(\alpha_0, \beta_0)$ is updated by observed
successes and failures to yield the posterior
$\text{Beta}(\alpha_0 + s, \beta_0 + f)$, where $s$ and $f$ count successes
and failures respectively \parencite{degroot1970}.

The well-known problem with Bayesian updating in \emph{non-stationary}
environments is that the posterior becomes increasingly concentrated around the
historical frequency as more data accumulates, making it progressively harder
to shift beliefs in response to a genuine structural change. We term this the
\emph{Sticky Prior Problem}: the weight of past evidence overwhelms the
evidence from a regime change, causing the posterior mean to ``stick'' near its
pre-change value for $O(\sqrt{n})$ additional observations. This phenomenon
is related to the concept of \emph{Bayesian regret} in non-stationary bandit
problems studied by \textcite{besbes2014stochastic}.

%---------------------------------------------------------------
\section{Multi-Armed Bandits: Exploration vs.\ Exploitation}
%---------------------------------------------------------------

The Multi-Armed Bandit (MAB) problem, formalized by \textcite{robbins1952},
provides the canonical framework for sequential decision-making under
uncertainty. An agent faces $K$ arms with unknown reward distributions and
must decide at each step which arm to pull so as to maximize cumulative
reward---balancing \emph{exploitation} of the currently best-known arm with
\emph{exploration} of potentially better alternatives.

\subsection{Upper Confidence Bound (UCB) Algorithms}

\textcite{auer2002finite} proved that the UCB1 algorithm achieves
$O(\log T)$ cumulative regret in the stochastic MAB setting. The algorithm
selects the arm $i$ with the highest index
\[
    \bar{\mu}_i(t) + \sqrt{\frac{2\ln t}{n_i(t)}},
\]
where $\bar{\mu}_i(t)$ is the empirical mean reward of arm $i$ and $n_i(t)$
is the number of times it has been pulled. The exploration bonus
$\sqrt{2\ln t / n_i}$ decays as an arm is pulled more frequently,
automatically balancing exploration and exploitation. This logarithmic regret
bound is known to be minimax optimal \parencite{lattimore2020bandit}.

\subsection{Thompson Sampling}

Thompson Sampling, introduced by \textcite{thompson1933likelihood}, maintains
a Bayesian posterior over arm reward probabilities and at each step samples
one parameter vector from the posterior, then selects the arm that appears
best under the sampled parameter. Despite its simplicity, Thompson Sampling
achieves $O(\log T)$ regret with favourable empirical constants
\parencite{agrawal2012analysis, russo2018tutorial}. Its Bayesian nature makes
it naturally suited to the Kelly betting context: the sampled probability
$\tilde{p}$ is used directly in the Kelly formula $f^*(\tilde{p})$, which
automatically reduces leverage when posterior uncertainty is high.

\subsection{EXP3 for Adversarial Environments}

When the reward sequence can be chosen adversarially---with no stochastic
structure---the standard UCB and Thompson Sampling regret bounds no longer
hold. \textcite{auer2002nonstochastic} introduced EXP3 (Exponential-weight
algorithm for Exploration and Exploitation), which maintains a probability
distribution over arms updated by exponential re-weighting:
\[
    w_i(t+1) = w_i(t) \cdot \exp\!\Bigl(\frac{\gamma \hat{r}_i(t)}{K}\Bigr),
\]
where $\hat{r}_i(t) = r_i(t)/p_i(t)$ is an importance-weighted reward
estimate and $\gamma$ is an exploration parameter. EXP3 achieves
$O(\sqrt{KT\ln K})$ regret against the best fixed arm in hindsight, with no
assumption on the reward generating process. Our implementation adapts EXP3
to continuous fractional betting; deviations from the canonical algorithm
are discussed in Section~\ref{sec:agents}.

%---------------------------------------------------------------
\section{Hidden Markov Models for Regime Detection}
%---------------------------------------------------------------

Hidden Markov Models (HMMs) were introduced to the statistics literature by
\textcite{baum1966statistical} and brought to prominence in time-series
econometrics by \textcite{hamilton1989new}, who used a two-state HMM to
model business cycle regime switching in GDP growth. The model posits that
an observed sequence $\{r_t\}$ is driven by an unobserved Markov chain
$\{S_t\}$ with transition matrix $A$, where each state $S_t = j$ emits
observations according to a regime-specific distribution $b_j(\cdot)$.

Inference over the hidden state is performed by the forward algorithm, which
computes the filtered probability
\[
    \alpha_t(j) = \Prob(S_t = j \mid r_1, \ldots, r_t),
\]
using the recursion
\[
    \alpha_t(j) = b_j(r_t) \sum_{i} \alpha_{t-1}(i) A_{ij}.
\]
We implement this recursion in log-space (using \texttt{logsumexp}) to
prevent numerical underflow over long sequences---a critical engineering
concern for $T > 250$ steps with float64 arithmetic.

A key limitation of standard HMMs is precisely the sticky prior problem
identified above: if the transition matrix is estimated from historical data
with high persistence (e.g., $A_{ii} \approx 0.99$), the model is slow to
assign probability to a regime change even when new evidence strongly suggests
one has occurred. Our Vol-HMM addresses this by (i) capping $A_{ii} \leq
0.95$ and (ii) augmenting the scalar return observation with a rolling
volatility feature, increasing the discriminative power of each observation.

The use of volatility as a regime indicator has empirical support from the
financial literature. \textcite{ang2002regime} show that the equity
risk premium is strongly counter-cyclical and that high-volatility regimes
predict low subsequent returns, validating the informational content of
volatility as a regime signal.

%---------------------------------------------------------------
\section{CUSUM Change-Point Detection}
%---------------------------------------------------------------

The Cumulative Sum (CUSUM) procedure, introduced by \textcite{page1954},
is a classical sequential test for detecting a shift in the mean of a
process. Given standardised observations $z_t = (r_t - \mu_0)/\sigma_0$,
the one-sided CUSUM statistic is
\[
    S_t^- = \max(0,\, S_{t-1}^- - z_t - k),
\]
where $k > 0$ is a drift parameter. An alarm is raised when $S_t^- >
h$ for some threshold $h$. CUSUM has optimal average run length properties
among all sequential tests for detecting a mean shift of a specified
magnitude \parencite{moustakides1986optimal}. We combine CUSUM with the HMM
forward algorithm to create a hybrid detector: CUSUM provides a rapid,
distribution-free alarm, while the HMM provides probabilistic regime
attribution.

%---------------------------------------------------------------
\section{Portfolio Insurance and CPPI}
%---------------------------------------------------------------

Constant Proportion Portfolio Insurance (CPPI) was introduced by
\textcite{black1987simplifying} as a dynamic strategy that maintains a floor
on wealth. The agent computes a ``cushion'' $C_t = W_t - F_t$ where $F_t$ is
a floor (often set as a fraction of peak wealth), and invests $m \cdot C_t$
in the risky asset, where $m > 1$ is a multiplier. In continuous time, CPPI
guarantees $W_t \geq F_t$ for all $t$; in discrete time, a single large
negative return can cause the cushion to turn negative (a ``gap risk'').

\textcite{busseti2016risk} extend the CPPI framework to the Kelly setting,
showing how a risk budget constraint on maximum drawdown can be incorporated
into the Kelly optimisation problem as a convex constraint. Our
\texttt{RiskConstrainedKelly} implementation follows their architecture,
using a trailing peak drawdown constraint $D_t = (W_{\text{peak}} -
W_t)/W_{\text{peak}} \leq D_{\max}$ with CPPI-style leverage
$\lambda_t = \min(1, m \cdot C_t / W_t)$.

%---------------------------------------------------------------
\section{Heavy Tails in Financial Returns}
%---------------------------------------------------------------

Empirical stock returns are well-documented to exhibit excess kurtosis
(``fat tails'') inconsistent with the Gaussian distribution assumed by
standard financial models. \textcite{mandelbrot1963variation} first
documented this for cotton prices, and the result has been replicated
extensively \parencite{cont2001empirical}. Practically, this means that
extreme events (e.g., daily moves of 5--10 standard deviations) occur far
more frequently than Gaussian models predict.

For the Kelly criterion, heavy tails are particularly dangerous: if an agent
sizes its position based on a Gaussian estimate of $\sigma$ but the true
distribution has heavier tails, there is a nonzero probability of a single
return $r_t$ satisfying $1 + f r_t \leq 0$---i.e., instantaneous ruin. Our
Experiment 6 (Section~\ref{ch:experiments}) demonstrates this
\emph{Kelly-Ruin Paradox} computationally by comparing ruin rates under
Gaussian versus Student-$t(\nu=3)$ return models using an identical Kelly
fraction calibrated to the Gaussian.

%---------------------------------------------------------------
\section{Summary of Identified Gaps}
%---------------------------------------------------------------

The literature reveals three specific gaps that this thesis addresses:

\begin{enumerate}
    \item \textbf{No filtration-consistent Kelly MAB framework.} Existing
    bandit algorithms are not designed with the log-utility objective, and
    existing Kelly implementations do not handle the MAB exploration problem.
    We bridge this gap via the Thompson-Kelly and UCB-Kelly agents.

    \item \textbf{Standard HMMs are too slow for adversarial regime changes.}
    The high persistence typically estimated in HMM transition matrices causes
    the sticky prior problem. We address this with the Vol-HMM's constrained
    transition matrix and CUSUM augmentation.

    \item \textbf{Kelly theory ignores finite-time ruin under fat tails.}
    Asymptotic optimality is irrelevant if the agent is ruined before the law
    of large numbers kicks in. We address this with the CPPI drawdown floor,
    quantifying the cost of this protection in terms of foregone growth.
\end{enumerate}
